\begin{thebibliography}{99}
\addcontentsline{toc}{section}{Daftar pustaka}
\bibitem{bahmani}
Bahmani-Oskooee, M. dan Fariditavana, H. (2016; terbit daring 2015).
Nonlinear ARDL Approach and the J-Curve Phenomenon.
\emph{Open Economies Review}, 27(1), 51--70.
\url{https://doi.org/10.1007/s11079-015-9369-5}.
Paper lengkap yang diberikan pengguna tersedia pada referensi proyek.

\bibitem{pss}
Pesaran, M. H., Shin, Y. dan Smith, R. J. (2001).
Bounds testing approaches to the analysis of level relationships.
\emph{Journal of Applied Econometrics}, 16(3), 289--326.
\url{https://doi.org/10.1002/jae.616}.

\bibitem{shin}
Shin, Y., Yu, B. dan Greenwood-Nimmo, M. (2014).
Modelling Asymmetric Cointegration and Dynamic Multipliers in a Nonlinear ARDL Framework.
Dalam \emph{Festschrift in Honor of Peter Schmidt}, 281--314.
Springer. \url{https://doi.org/10.1007/978-1-4899-8008-3_9}.

\bibitem{ks}
Kripfganz, S. dan Schneider, D. C. (2023).
ardl: Estimating autoregressive distributed lag and equilibrium correction models.
\emph{The Stata Journal}, 23(4), 983--1019.
\url{https://doi.org/10.1177/1536867X231212434}.

\bibitem{bh}
Benjamini, Y. dan Hochberg, Y. (1995).
Controlling the False Discovery Rate: A Practical and Powerful Approach to Multiple Testing.
\emph{Journal of the Royal Statistical Society: Series B}, 57(1), 289--300.
\url{https://doi.org/10.1111/j.2517-6161.1995.tb02031.x}.

\bibitem{bec}
United Nations Statistics Division (2018).
\emph{Classification by Broad Economic Categories, Revision 5},
ST/ESA/STAT/SER.M/53/Rev.5.
\url{https://unstats.un.org/unsd/trade/classifications/bec.asp}.
Manual dan korespondensi disimpan dalam referensi proyek.

\bibitem{unsdcorr}
United Nations Statistics Division (2022).
\emph{HS--SITC--BEC Correlations dan korespondensi antar revisi HS}.
\url{https://unstats.un.org/unsd/classifications/Econ/}.
File sumber lokal: \texttt{unsd\_correlations\_2022.xlsx}
dan korespondensi HS--SITC dalam folder referensi klasifikasi.

\bibitem{unctad}
UNCTAD (tanpa tahun; salinan lokal diakses 2026).
\emph{SITC Rev.3 Products: Special Groupings Hierarchy},
kelompok komoditas A11.
\url{https://unctadstat.unctad.org/EN/Classifications/DimSitcRev3Products_DsibSpecialGroupings_Hierarchy.pdf}.
Salinan lokal dan manifest unduhan dipertahankan untuk penelusuran.

\bibitem{stata}
StataCorp.
\emph{estat sbcusum: Cumulative sum test for parameter stability}.
Manual resmi Stata.
\url{https://www.stata.com/manuals/tsestatsbcusum.pdf}.

\bibitem{local}
Dokumentasi proyek Tio (2026).
Audit data HS-6, audit validitas makro, manifest sumber dan hasil estimasi.
Berkas lokal pada \path{input/reference/},
\path{output/diagnostics/data/} dan
\path{output/diagnostics/oskooee/}.
Audit cakupan negara pelapor: \path{output/presentation/scope_audit_indonesia.json}.
Statistik dan estimasi naskah menggunakan keluaran run 17 November 2026.
\end{thebibliography}
